\section{Negative-weight background subtraction}
\label{app:negative-weight-background}
\label{app:negweight}

\subsection{Estimator}
\label{sec:negweight-estimator}

Negative-weight subtraction supplies the reconstructed signal target for the first
OmniFold step without imposing a subtraction grid. Using the notation of
Sec.~\ref{sec:omnifold-implementation}, we give data unit weight and append
simulated background with the negative of its exposure-scaled weights. The
measured class then has signed density
\begin{equation}
  \rho_1^{\mathrm{neg}}(x)=D(x)-B(x)
  \label{eq:negweight-continuous}
\end{equation}
and the simulated class has the current selected-signal density $S(x)$.
For $S(x)>0$ and $D(x)>B(x)$, the population classifier odds with unnormalized
class weights are $[D(x)-B(x)]/S(x)$. In a bin with positive net yield, purity
weighting instead scales the data shape by the bin-averaged signal fraction.
\begin{equation}
  \rho_1^{\mathrm{pur}}(x)=D(x)\,w_{\mathrm{pur}}(b)
  =D(x)\,\max\!\left(0,1-\frac{B_b}{D_b}\right)
  \label{eq:negweight-purity}
\end{equation}
The signed target differs from the purity-weighted target by
$D(x)B_b/D_b-B(x)$, with zero net change in the bin.

The refinement gives the first-step classifier a nonnegative target. The
sign classifier is trained on the absolute-weight sample, with density
$D+B$. Its population score is
$g=D/(D+B)$ wherever $D+B>0$. The factor $\max(0,2g-1)$ then gives
the target $[D-B]_+$, where $[z]_+\equiv\max(0,z)$~\cite{Nachman:2025staypositive}.
For a finite learner, the clipping acts on the fitted score. With
histogram-based boosting in scikit-learn, the direct signed fit and the
clipped refinement gave a median cross-section bin ratio of $\nwMedianBin$. With
exact-split trees, the direct fit produced large excursions at low $p_T$
that the refinement removed.

\subsection{Binned limit}
\label{sec:negweight-binned}

The signed and purity-weighted targets give identical updates where the
classifier is constant within bins and the net bin yields are nonnegative.
For an occupied reconstructed bin $b$ on common support, with simulated
signal yield $S_b>0$, the signed binary cross-entropy contribution is
\begin{equation}
  \mathcal L_b(p_b)=-(D_b-B_b)\log p_b-S_b\log(1-p_b).
  \label{eq:background-binned-loss}
\end{equation}
For $D_b>B_b$, minimizing this loss gives
\begingroup
\refstepcounter{equation}
\label{eq:background-binned-ratio}
\label{eq:negweight-rb}
\begin{equation*}
  r_b=\frac{p_b}{1-p_b}=\frac{D_b-B_b}{S_b}.
  \tag{\theequation}
\end{equation*}
\endgroup
Purity weighting supplies the same class-1 mass by
Eq.~\eqref{eq:purity-weight}. At $D_b=B_b$, the limiting optimum has $p_b=0$
and $r_b=0$. The correspondence between binned OmniFold and iterative
Bayesian unfolding makes Eq.~\eqref{eq:background-binned-ratio} the usual
subtract-then-unfold update~\cite{Andreassen:2019cjw,Andreassen:2021zzk,DAgostini:1994fjx}.

For $D_b<B_b$, the signed loss tends to $-\infty$ as $p_b\to0^+$, so setting the
yield to zero is an added prescription. If the sign classifier is itself
restricted to the bins, the refinement gives the same mass $[D_b-B_b]_+$ as
purity weighting. An unbinned refinement clips before integration, so its
population bin mass can exceed $[D_b-B_b]_+$ if the net density has both signs
within the bin.

\subsection{Comparison with purity subtraction}
\label{sec:negweight-2d}
\label{sec:negweight-uq}
\label{sec:negweight-footing}

The two treatments give similar integrated cross sections in two dimensions.
Using the scikit-learn histogram-based learner and a common training seed,
purity and negative-weight subtraction give totals of
\num[exponent-mode=fixed,fixed-exponent=-38,drop-exponent=true]{\nwSigPur} and
\num[exponent-mode=fixed,fixed-exponent=-38,drop-exponent=true]{\nwSigNeg},
respectively, in units of $10^{-38}\,\mathrm{cm}^2$ per nucleon, a change of
\SI{\nwPctTot}{\percent}. The negative-weight to purity ratio across populated
cross-section bins has a median of $\nwMedianBin$ and a root-mean-square
deviation from unity of \SI{\nwRmsBin}{\percent}. The largest shift is
\SI{\nwWorstBin}{\percent} in a low-$p_T$ edge bin with a large background
contribution (Table~\ref{tab:negweight} and Fig.~\ref{fig:negweight-ratio}).

\begin{table}[!hbp]
  \centering
  \caption{Purity versus negative-weight background subtraction in 2D}
  \label{tab:negweight}
  \begin{tabular}{lcc}
    \hline\hline
    & binned purity & negative-weight \\
    \hline
    total $\sigma$ (cm$^2$/nucleon) & $\num{\nwSigPur}$ & $\num{\nwSigNeg}$ \\
    total ratio (neg/pur)           & \multicolumn{2}{c}{$\nwRatioTot$ (\SI{\nwPctTot}{\percent})} \\
    median per-bin ratio            & \multicolumn{2}{c}{$\nwMedianBin$} \\
    per-bin RMS deviation           & \multicolumn{2}{c}{\SI{\nwRmsBin}{\percent}} \\
    largest single-bin deviation    & \multicolumn{2}{c}{\SI{\nwWorstBin}{\percent} (low $p_T$, high background)} \\
    \hline\hline
  \end{tabular}
\end{table}

\begin{figure}[!hbp]
  \centering
  \includegraphics[width=0.7\textwidth]{negweight_ratio_2d}
  \caption{The ratio of the negative-weight to the binned-purity cross section has a
  median of $\nwMedianBin$, and only the low-$p_T$, high-background edge bin
  departs by more than a few percent. The map shows the ratio in the populated
  bins of the 2D measurement, with a common estimator and seed. The color scale
  saturates.}
  \label{fig:negweight-ratio}
\end{figure}

The refined estimator gives a smaller nominal-truth closure residual in the
bins most sensitive to subtraction. In the 20 paired background-inclusive
pseudo-experiments of Sec.~\ref{sec:muon-cross-section-validation}, the largest
difference between the two mean estimates is $5.45\sigma_{\mathrm{stat}}$ in
the bin $0.70\le p_T<0.85\,\mathrm{GeV}/c$,
$1.5\le p_\parallel<2.0\,\mathrm{GeV}/c$. Here $\sigma_{\mathrm{stat}}$ is
from the purity-based statistical covariance and serves as a common reference
scale for the descriptive comparison.

On data, switching from purity to the refined estimator shifts the integral
by $-0.139\%$. The median absolute bin shift is $0.44\%$ and the maximum is
$10.1\%$. The shifts exceed $\sigma_{\mathrm{stat}}$ in 43 of 205 bins and
are below half the purity-based total standard deviation in every bin.
Table~\ref{tab:estimators} gives the subtraction treatment for each
reported result.

Both systematic and statistical variations enter before the background
correction. A systematic change to the simulated background enters through
its varied weights or reconstructed selection, as applicable. In the
statistical bootstrap, independent Poisson multiplicities for data, signal
simulation, and background simulation enter each refit. Each unfolding uses
the background correction formed from its own inputs, including a new sign
fit when using the refined estimator. Appendix~\ref{app:covariance-construction}
gives the ensemble definitions.

In a matched, full-statistics LightGBM check in two dimensions, the ratios
of $\sqrt{\operatorname{tr} C}$ for negative-weight to purity subtraction are
$\nwSystRatio$ for the 187-universe systematic ensemble and $\nwStatRatio$ for
the statistical ensemble with matched training seeds. The trace ratios
quantify agreement only in the aggregate diagonal variance.

A common shift from changing the subtraction cancels in a mean-centered
covariance. If the change shifts the cross-section vector of member $u$ by
$\delta_u$, the centered vector shifts by $\delta_u-\bar\delta$, where
$\bar\delta$ is the mean shift within the systematic band or the bootstrap
ensemble.

The five-dimensional systematic check covers two individual variations at
one unfolding iteration with the scikit-learn histogram-based learner. For the
two-particle--two-hole variation, both treatments give an integrated cross
section of $3.014\times10^{-38}\,\mathrm{cm}^2$ per nucleon at the reported
precision. For the quasielastic axial-mass variation, the integral shifts by
$-0.03\%$. The respective median bin ratios are 1.0005 and 0.9990. The
five-dimensional background-inclusive closure tests probe bias separately
(Sec.~\ref{sec:higher-dimensional-closure}).

\subsection{Background treatment in other unbinned analyses}
\label{sec:negweight-practice}

Unbinned measurements have handled backgrounds through event selection,
event weights, and systematic variations. H1 suppressed backgrounds to a
negligible level with its selection, and LHCb applied track-by-track
corrections before the unbinned detector unfolding~\cite{H1:2021wkz,LHCb:2022rky}.
In its 24-observable $Z+$jets measurement, ATLAS assigned an uncertainty for
the residual background, which was below
$0.25\%$. The analysis and subsequent guidance describe adding negatively
weighted simulated events to the measured class for larger
backgrounds~\cite{ATLAS:2024jets,Canelli:2025guide}.
In the full charged-particle phase-space measurement, ATLAS used negatively
weighted background simulation to train a neural network whose outputs
reweight data before unfolding~\cite{ATLAS:2026fullps}.
OmniFold-HI models large backgrounds in the unfolding, while generative
approaches can subtract backgrounds before correcting for acceptance and
inefficiency~\cite{Falcao:2025backgrounds,Butter:2025genfold}.

% PASS2:
% No macros in the source for the 20 pseudo-experiments, 5.45-sigma shift,
% bin edges (0.70, 0.85, 1.5, 2.0), data shifts (-0.139%, 0.44%, 10.1%),
% 43 of 205 bins, 187 universes, five-dimensional checks (3.014 x 10^-38,
% -0.03%, 1.0005, 0.9990), or 0.25% ATLAS background. These are source literals.
% "Half" paraphrases the source's literal 0.497. The ATLAS count 24 is
% spelled out in the original introduction and bibliography, with no macro.
% The two totals print the macro mantissas in the prose's common 10^-38 units.
% All local labels and reference targets are accounted for. The former
% propagation and production subsections share the merged comparison number.
% The restored table uses the source's short caption as its visible caption.
